\section{1. Información general del
proyecto}\label{informaciuxf3n-general-del-proyecto}

\subsection{1.1 Integrantes y roles}\label{integrantes-y-roles}

{\def\LTcaptype{none} % do not increment counter
\begin{longtable}[]{@{}
  >{\raggedright\arraybackslash}p{(\linewidth - 8\tabcolsep) * \real{0.2456}}
  >{\raggedright\arraybackslash}p{(\linewidth - 8\tabcolsep) * \real{0.1282}}
  >{\raggedright\arraybackslash}p{(\linewidth - 8\tabcolsep) * \real{0.2012}}
  >{\raggedright\arraybackslash}p{(\linewidth - 8\tabcolsep) * \real{0.1686}}
  >{\raggedright\arraybackslash}p{(\linewidth - 8\tabcolsep) * \real{0.2564}}@{}}
\toprule\noalign{}
\begin{minipage}[b]{\linewidth}\raggedright
\textbf{Nombre completo}
\end{minipage} & \begin{minipage}[b]{\linewidth}\raggedright
\textbf{Matrícula}
\end{minipage} & \begin{minipage}[b]{\linewidth}\raggedright
\textbf{Usuario GitHub}
\end{minipage} & \begin{minipage}[b]{\linewidth}\raggedright
\textbf{Rol(es) y sprints}
\end{minipage} & \begin{minipage}[b]{\linewidth}\raggedright
\textbf{Áreas principales de contribución}
\end{minipage} \\
\midrule\noalign{}
\endhead
\bottomrule\noalign{}
\endlastfoot
Abdías Daniel Chávez Cruz & 0224170024 & AbdiasChavez & Product Owner
(S0--S6)

Developer & Gestión del Product Backlog, relación con el cliente,
diagrama del proceso actual (3.5), base de datos. \\
Ángel Yael Avendaño Guzmán & 0224170023 & AyagDan & Scrum Master
(S0--S6)

Developer & Administración del repositorio, revisión e integración de
Pull Requests, documento inicial. \\
Dulce María Estudillo Rosado & 0224170028 & DulceRosado & Developer
(S0--S6) & Requisitos funcionales y no funcionales (4.4, 4.6),
stakeholders (3.4). \\
Isis Alhelí González Maldonado & {[}Pendiente{]} &
IsisAlheliGonzalezMaldonado & Developer (S0--S6) & Historias de usuario
y criterios de aceptación (4.5). \\
Isis Yamilet Meléndez Hernández & 0224170037 & IsisMelendezHernandez &
Developer (S0--S6) & Integración del machote, descripción general (4.3),
reglas de negocio (4.7), protección de datos (4.8). \\
\end{longtable}
}

\subsection{1.2 Enlaces del proyecto}\label{enlaces-del-proyecto}

{\def\LTcaptype{none} % do not increment counter
\begin{longtable}[]{@{}
  >{\raggedright\arraybackslash}p{(\linewidth - 2\tabcolsep) * \real{0.3631}}
  >{\raggedright\arraybackslash}p{(\linewidth - 2\tabcolsep) * \real{0.6369}}@{}}
\toprule\noalign{}
\begin{minipage}[b]{\linewidth}\raggedright
\textbf{Recurso}
\end{minipage} & \begin{minipage}[b]{\linewidth}\raggedright
\textbf{Enlace}
\end{minipage} \\
\midrule\noalign{}
\endhead
\bottomrule\noalign{}
\endlastfoot
Repositorio & \url{https://github.com/AbdiasChavez/Florentino} (original:
\url{https://github.com/AyagDan/SistemaDeServiciosYCitas}) \\
Integración continua & \url{https://github.com/AbdiasChavez/Florentino/actions} \\
Tablero (GitHub Projects / Jira) & \pendiente{crear el tablero en GitHub Projects con las HU-01 a
HU-30 como issues}. \\
Prototipo de UI (Figma / Penpot) & \pendiente{se agrega al crear el prototipo (Sprint 4)}. \\
Ambiente de pruebas (staging) & No aplica todavía (desde el Sprint
4). \\
Release v0.1.0 & No aplica todavía (Demo final). \\
Canal de comunicación del equipo & Pendiente: registrar el canal
acordado (ver CONTRIBUTING.md). \\
\end{longtable}
}

\subsection{1.3 Resumen ejecutivo}\label{resumen-ejecutivo}

Servicio Automotriz Hermanos Ruiz opera hoy con cuadernos, órdenes en
papel y llamadas. Esto provoca disputas por cobros que el cliente dice
no haber autorizado, reclamos por daños sin evidencia del estado en que
llegó el vehículo, muchas llamadas para preguntar por el avance y poca
visibilidad de qué vehículo ocupa cada una de las 4 bahías.

El Sistema de Órdenes de Servicio y Citas digitaliza el ciclo de
atención: la cita, la recepción del vehículo con checklist y
fotografías, el diagnóstico y la asignación de bahía, la cotización que
el cliente autoriza desde su celular con un código de verificación, y el
cobro y la entrega. El cliente consulta el estado de su vehículo con un
enlace seguro, sin instalar aplicaciones.

Estado actual (v\versiondoc): se levantaron los requisitos con el
cliente (entrevista del 19/10/2026) y se definió el alcance. Ya funciona,
en una versión de desarrollo sin pantallas todavía, el corazón del
sistema: abrir una orden con checklist y fotografías, asignar bahía y
mecánico, registrar el diagnóstico, enviar la cotización al cliente,
autorizarla con el código de verificación, cobrar y entregar sin saldo
pendiente. Cada una de estas reglas se comprueba automáticamente con 35
pruebas. Falta: las pantallas para el personal y para el cliente, el envío
real de correos, la operación sin internet, el almacén y el ambiente de
pruebas en línea. La Fase 2 agregará reportes gerenciales, consulta de
existencias y de garantías, y el tablero simplificado de bahías.
